\section{Introduction}\label{p2:sec-intro}

Points of $\R^{2n}=\C^n$ are $z=(z_1,\dots,z_n)$, $z_j=x_j+iy_j$, with $\omega_0=\sum_jdx_j\wedge dy_j$. The
Hamiltonian vector field of $H(t,z)$ is defined by $\iota_{X_H}\omega_0=dH$, i.e.\
$X_H=(\partial_yH,-\partial_xH)$, and $\phi^t_H$ is its flow from time $0$ to time $t$. We write
$B(c)=\{\pi|z|^2\le c\}$ for the closed ball of capacity $c$, $\Bo(c)$ for its interior, and $B(s,c)=s+B(c)$.

Balls of capacities $c_1,\dots,c_k$ can be packed symplectically into $\Bo(C)$ only if $\sum_ic_i^n<C^n$ (volume) and
$c_i+c_j<C$ for $i\ne j$ (Gromov's two-ball theorem \cite{Gromov1985}). In dimension four there are further
obstructions \cite{McDuffPolterovich1994,Hutchings2011}, and every packing problem reduces to inequalities indexed by
exceptional classes \cite{McDuff1998,LiLiu2001,LiLi2002}; Biran \cite{Biran1997} proved that $k\ge9$ equal balls fill
$B^4$ up to any $\eps>0$. In dimension $2n\ge6$ there are no further obstructions: this is Conjecture~A of Siegel and
Yao \cite{SiegelYao2025}, proved in a 2026 preprint of OpenAI \cite{OpenAI2026}. These proofs do not write down a
Hamiltonian moving prescribed balls; in \cite{OpenAI2026} the comparison function is a fixed point, and the planar maps
realising the rearrangements and the relative Moser isotopies are obtained from existence statements. Paper~I \cite{PaperI} constructs explicit Hamiltonian full packings of $B^4$ by $k\ge10$ equal
balls. Here we realise, by the time-one map of one explicit Hamiltonian moving balls at prescribed positions into the
target, every packing allowed by the two inequalities in dimension $2n\ge6$, and in dimension four a range that
contains all nearly full packings by $k\ge9$ equal balls; and we give an algorithm whose output is a rational
polynomial times a fixed cutoff.

\subsection{Results}\label{p2:sec-intro-results}

\begin{definition}[Packing specification]\label{def:intro-spec}
\lean{SymplecticPacking.Problem.PackingSpec}
Let $n\ge1$, $k\ge0$, $C>0$, $c=(c_1,\dots,c_k)\in(0,\infty)^k$ and $s=(s_1,\dots,s_k)\in(\R^{2n})^k$. A
function $H\colon\R\times\R^{2n}\to\R$ satisfies $\Pack(n,k,C,c,s)$ if
\begin{enumerate}[label=(\alph*)]
\item $H$ is $C^\infty$, and there is a compact set $\mathcal S\subset\R^{2n}$ with $H(t,z)=0$ for all
 $t\in\R$ and all $z\notin\mathcal S$;
\item for every $i$ and every $z\in B(s_i,c_i)$ there is an integral curve $\gamma\colon\R\to\R^{2n}$ of
 $X_H$ with $\gamma(0)=z$, and every integral curve $\gamma\colon[0,1]\to\R^{2n}$ of $X_H$ with
 $\gamma(0)=z$ satisfies $\gamma(1)\in\Bo(C)$.
\end{enumerate}
\leanmark{Problem.PackingSpec}
\end{definition}

Under (a) the flow $\phi^t_H$ is global and symplectic, and (b) says $\phi^1_H(B(s_i,c_i))\subset\Bo(C)$.

\begin{definition}[Admissible data]\label{def:intro-admissible}
\lean{SymplecticPacking.Problem.Admissible}
Let $n\ge2$, $k\ge1$, $C>0$ and $c\in(0,\infty)^k$. The data $(n,k,C,c)$ are \emph{admissible} if
\begin{align*}
 n\ge3:&\qquad \textstyle\sum_ic_i^n<C^n\quad\text{and}\quad c_i+c_j<C\ \text{ for all } i\ne j;\\
 n=2:&\qquad \textstyle\sum_ic_i^2<C^2\quad\text{and}\quad 3c_i<C\ \text{ for all } i.
\end{align*}
\leanmark{Problem.Admissible}
\end{definition}

For $n\ge3$ these are exactly the two necessary conditions. For $n=2$ they define a sufficient range: by \cite{McDuffPolterovich1994,McDuff1998,LiLiu2001,LiLi2002}, closed balls of capacities
$c_i$ embed in $\Bo(C)$ if $\sum_ic_i^2<C^2$ and $\sum_im_ic_i<dC$ for every exceptional class $(d;m_1,\dots,m_k)$
with $d\ge1$; such a class has $m_i\ge0$ and $\sum_im_i=3d-1$, so $3c_i<C$ gives $\sum_im_ic_i<dC$.


\begin{maintheorem}\label{thm:main}
\uses{def:intro-spec, def:intro-admissible, def:intro-levels, def:ham-H-final, thm:eff-algorithm}
Let $n\ge2$, $k\ge1$, $C>0$, $c_1,\dots,c_k>0$ and $s_1,\dots,s_k\in\R^{2n}$ be such that the closed balls
$B(s_i,c_i)$ are pairwise disjoint and $(n,k,C,c)$ is admissible. Then:
\begin{enumerate}[label=(\roman*)]
\item the function $H=H[n,k,C,c,s]$ of Definition~\ref{def:ham-H-final} is $C^\infty$ on $\R\times\R^{2n}$,
 vanishes outside $\R\times\mathcal S$ for an explicit compact set $\mathcal S$ and for $t\notin(0,1)$, and
 satisfies $\Pack(n,k,C,c,s)$; in particular $\phi^1_H$ maps each $B(s_i,c_i)$ into $\Bo(C)$, and the
 images are pairwise disjoint;
 \leanmark{main\_theorem, main\_time\_support, main\_support\_radius, flow\_images\_disjoint}
\item \textup{(L1)} the definition of $H[n,k,C,c,s]$ is the same for all such data and uses only the
 operations of level~L1 in Definition~\ref{def:intro-levels};
\item \textup{(L2, L3)} there is one oracle algorithm which, given $n$, $k$ and names of the real inputs
 $C,c_1,\dots,c_k,s_1,\dots,s_k$ (a rational number is named by its constant sequence), halts on all such data and
 outputs an integer $B$ and a polynomial $P\in\Q[t,z]$ such that $\Hfin=\beta_B\cdot P$ satisfies
 $\Pack(n,k,C,c,s)$. The output may depend on the names; for computable inputs, given by computable names, it is an
 ordinary algorithm.
 \leanmark{original\_name\_algorithm}
\end{enumerate}
\end{maintheorem}

Here level L1 means that $H$ is built by explicitly listed operations (arithmetic, elementary functions, integrals,
unique solutions of explicit ODEs, inverses, case distinctions, finite searches; Definition~\ref{def:intro-levels}).
A \emph{name} of a real number $x$ is a sequence of rationals $q(p)$ with $|q(p)-x|\le2^{-p}$, and
$\beta_B(z)=\prod_l\beta(\zeta_l/B)$ is a fixed smooth cutoff equal to one on the cube $[-B/2,B/2]^{2n}$ and vanishing outside
$(-3B/4,3B/4)^{2n}$, so that L2 is a finite description of a Hamiltonian and L3 a uniform algorithm producing it.

Parts (i) and (ii) are proved in \S\ref{p2:sec-proof}, from the model theorems of Section~\ref{p2:sec-models} and the
passage from embeddings to one Hamiltonian of Paper~I; for $k\ge10$ equal capacities in dimension four with
$\frac34<\rho\sqrt k<1$, $\rho=c/C$, they are the linear branch of Paper~I. Part (iii) is
Theorem~\ref{thm:eff-algorithm}, proved in Section~\ref{sec:eff}; its termination argument uses part (i).

\begin{maincorollary}[The Siegel--Yao criterion by explicit flows]\label{thm:main-sy}
Let $n\ge3$, $k\ge1$, $C>0$ and $c_i>0$, and let $B(s_i,c_i)\subset\R^{2n}$, $1\le i\le k$, be pairwise disjoint closed balls.
There is a symplectic embedding of a neighbourhood of $\bigsqcup_iB(s_i,c_i)$ into $\Bo(C)$ if and only if
$\sum_ic_i^n<C^n$ and $c_i+c_j<C$ for all $i\ne j$. In that case $\phi^1_H$ with $H=H[n,k,C,c,s]$ is such
an embedding.
\leanmark{neighbourhood\_criterion, flow\_neighbourhood\_embedding}
\end{maincorollary}

Sufficiency is the main theorem of \cite{OpenAI2026}; Corollary~\ref{thm:main-sy} realises the embedding as the
time-one map of an explicit Hamiltonian moving balls at prescribed positions.

\begin{maincorollary}[Equal balls]\label{thm:main-equal}
Let $n\ge2$, $\eps\in(0,1)$, and $k\ge2^n$ if $n\ge3$, $k\ge9$ if $n=2$. Let $c>0$, let $s_1,\dots,s_k$ be
centres with $B(s_i,c)$ pairwise disjoint, and let $C>0$ with $kc^n<C^n<kc^n/(1-\eps)$. Then
$H[n,k,C,(c,\dots,c),s]$ satisfies $\Pack(n,k,C,(c,\dots,c),s)$, and $kc^n/C^n>1-\eps$.
\leanmark{equal\_balls}
\end{maincorollary}

Thus, for $n\ge3$, $k\ge2^n$ equal balls fill any fraction below one of $B(C)$ by an explicit flow; for $2\le k<2^n$
the two-ball inequality caps the fraction at $k/2^n$. For $n=2$ the bound $k\ge9$ is Biran's \cite{Biran1997}; here it
comes from the condition $3c<C$, and for $k\le8$ full packings do not exist except for $k=1,4$
\cite{McDuffPolterovich1994}. The case $n=2$ gives a proof, without $J$-holomorphic curves or Seiberg--Witten theory, of Biran's
theorem that $B^4$ is fully packed by $k$ equal balls for every $k\ge9$; the earlier proofs \cite{Biran1997,Biran1999,BuseHind2013,BuseHindOpshtein2016}
use them, directly or through the results they cite.

\begin{maincorollary}[A cone of unequal balls in dimension four]\label{thm:main-cone4}
Let $k\ge1$ and $c_1,\dots,c_k>0$ with $\max_ic_i\le\frac13\big(\sum_ic_i^2\big)^{1/2}$, let $\eps\in(0,1)$ and
$C_\eps=\big(\sum_ic_i^2/(1-\eps/2)\big)^{1/2}$. For all centres $s_i$ with $B(s_i,c_i)$ pairwise disjoint,
$H[2,k,C_\eps,c,s]$ satisfies $\Pack(2,k,C_\eps,c,s)$, and $\sum_ic_i^2/C_\eps^2=1-\eps/2$; the images fill the fraction
$1-\eps/2$ of the volume of $B(C_\eps)$.
\leanmark{unequal\_cone}
\end{maincorollary}

\begin{maincorollary}[The Epoch data]\label{thm:main-epoch}
Let $k\ge9$ be an integer, $\eps\in(0,1)$, $s_i=(3i-3,0,0,0)\in\R^4$, $R=(k/(1-\eps/2))^{1/4}$ and $C=\pi R^2$. Then
$H=H[2,k,C,(\pi,\dots,\pi),s]$ satisfies $\Pack(2,k,C,(\pi,\dots,\pi),s)$: its time-one map sends each closed unit ball
$B(s_i,\pi)$ into $\Bo(\pi R^2)$, and $k/R^4=1-\eps/2>1-\eps$.
\leanmark{epoch\_packing}
\end{maincorollary}

For $k\ge10$ and $\eps<\frac78$ this is the corresponding corollary of Paper~I, which uses the linear comparison datum;
the other cases use the scalar comparison datum. The corollary is the input of the Epoch part of
Theorem~\ref{thm:eff-algorithm}, which, together with Paper~I, answers the mathematical part of the construction problem
of Epoch AI \cite{Epoch} (see Remark~\ref{rem:eff-epoch-interface} for its numerical interface).

\subsection{The construction and what is new}\label{p2:sec-intro-mechanism}

We follow the mechanism of \cite{OpenAI2026}, and its sign convention $\iota_{X_Q}\omega_V=-dQ$ for fibre
Hamiltonians, opposite to the convention for $H$. The model is a product $\Sigma_\ell\times\C^m$, $m=n-1$, of a compact
piece of an affine curve with the fibre, with fibre moments $p_j=\pi|w_j|^2$. A ball $B(a)\subset\C\times\C^m$,
sliced by the moments of its last $m$ coordinates, is a family of discs of capacity $(a-|p|)_+$, $|p|=\sum_jp_j$. For
capacities below $\frac12$ the base is the diagonal curve $y_j^2=t^2-j$, and the base area above a fibre point is
$\Hb(p)=A_\ell(1-2|p|)$ with $A_\ell\uparrow2^m$; the limiting profile has integral $1/n!=\vol\Bo(1)$. For a nearly
full packing the deficit $\Pb=\Hb-\sum_i(a_i-|p|)_+$ has positive integral but is negative near $p=0$, so the base
discs cannot simply be placed side by side. One twists the bundle across an annulus by a Hamiltonian isotopy $q$ of the
fibre; with a function $K$ of the moments this redistributes base area between moment levels, and the balls fit if
\[
 (K-P)\bigl(p(q(w))\bigr)\le K\bigl(p(w)\bigr)-\gamma\qquad\text{for every fibre point }w,
\]
for some $\gamma>0$ and a minorant $P\le\Pb$. In \cite{OpenAI2026}, $K$ is the fixed point of a contraction built
from increasing rearrangements, the rearrangements are realised by isotopies, the inequality is turned into
embeddings by Moser's method, and the result is transferred to the target. There are three cases: all capacities
below $\frac C2$; one capacity at least $\frac C2$ with $n\ge4$; and one capacity at least $\frac C2$ with $n=3$.

The common isotopy from the prescribed balls and the Hamiltonian generating it are those of Paper~I, recalled in
Section~\ref{p2:sec-machine}; so are the cubic model, the four-dimensional handle and the return to the four-ball.
The following ingredients are new here:
\begin{enumerate}[label=(\arabic*),leftmargin=*]
\item a comparison datum for every concave profile on a rational polytope: a finite rational iterate of the
 rearrangement contraction of \cite[\S2]{OpenAI2026}, with explicit gap, in place of its fixed point
 (Theorem~\ref{thm:cmp-main});
\item an explicit area-preserving rearrangement of a planar convex function, lifted to the fibres of a toric model
 with face collars (Theorem~\ref{thm:cmp-realisation}, proved in \Ssec{sec:plc});
\item in dimension four, a scalar comparison datum for the cubic model that covers all admissible capacities, including
 $k=9$ and unequal capacities, where a linear comparison function can fail (Theorem~\ref{thm:mdl-cub-general});
\item a bordered handle whose Moser maps are time-one maps of explicit fields in a finite coordinate atlas, which
 proves the packing theorem in a model (Theorem~\ref{thm:mech-packing}) for $n\ge3$
 (\Ssec{p2:sec-shd}--\Snum{sec:native-toric});
\item a model over the fixed diagonal curve $y_j^2=t^2-j$, an explicit instance of the complete-intersection model of
 \cite[\S5.2]{OpenAI2026}, with a transfer to the ball by explicit flows (Theorem~\ref{thm:mdl-qdr});
\item for one large ball and $n\ge4$, a mixed product model with a closed-form return to the shell about the reserved
 ball; for $n=3$, the polynomial cubic model of \cite{OpenAI2026} with a quantitative area tail and a polynomial
 bubble that inserts the large ball (Theorem~\ref{thm:mdl-lrg});
\item an algorithm with finite output: a sound rational Euler certificate for a candidate polynomial, and termination
 proved from the explicit Hamiltonian (Theorem~\ref{thm:eff-algorithm}).
\end{enumerate}
Definitions~\ref{def:intro-spec} and~\ref{def:intro-admissible}, Theorem~\ref{thm:main}(i), (iii) and
Corollaries~\ref{thm:main-sy}--\ref{thm:main-epoch} are verified in Lean~4; the explicitness statement (ii) is
checked by a metaprogram that is not itself verified (Section~\ref{p2:sec-ver}).

\subsection{Relation to previous work}\label{p2:sec-intro-prior}

Packing stability in dimension four was extended by Biran \cite{Biran1999}, Buse, Hind and Opshtein
\cite{BuseHind2011,BuseHind2013,BuseHindOpshtein2016} and Edtmair \cite{Edtmair2025}; Buse and Hind proved it in higher
dimensions. Siegel and Yao
\cite{SiegelYao2025} stated Conjecture~A and proved a stabilised two-ball theorem and the invariance of
four-dimensional ball packing problems under products with closed surfaces. Explicit constructions go back to
Traynor \cite{Traynor1995}, including full packings of $B^{2n}$ by $d^n$ equal balls; symplectic folding and wrapping
are in \cite{SchlenkHand,SchlenkBook,Schlenk2018}, and Witt Nystr\"om \cite{WittNystrom2024} embedded large
neighbourhoods of zero sections of normal bundles. It uses the rearrangement contraction of \cite[\S2]{OpenAI2026}, and the formal proof uses lemmas of the Lean library accompanying
\cite{OpenAI2026}; the packing existence theorem of that preprint is not used.

\subsection{Organisation}\label{p2:sec-intro-organisation}

Section~\ref{p2:sec-machine} states the comparison data and the toric mechanism, proved here, and the passage from
embeddings to one Hamiltonian, recalled from Paper~I. Section~\ref{p2:sec-models} describes the four-dimensional
scalar branch, the three higher-dimensional models and the case distinction. Section~\ref{p2:sec-proof} proves
Theorem~\ref{thm:main}(i), (ii) and the corollaries, Section~\ref{sec:eff} proves the finite form and the algorithm, and
Section~\ref{p2:sec-ver} describes the formal verification. The supplement contains the complete proofs; numbers with
the prefix S refer to it, and ``Paper~I, \dots'' refers to \cite{PaperI} and its supplement.
